\subsection{综合评价的排序可信度}
问题五的综合评价只用于三种情景之间的相对比较。四个指标经正向化后，熵权、CRITIC、等权和 2000 次局部扰动都给出同一排序，其中“禁渔+污染+入侵”始终最差，当前比较范围内，复合压力情景始终排在末位。

综合分数不能替代单指标解释。食源、江豚和长江鲟仍对应不同层面的生态状态；若后续补入独立生态监测量并扩展污染、入侵强度梯度，新增数据用于检查这一排序能否外推。

\noindent 沿“恢复收益--上传导--长期复杂化--复合压力检验”关系，问题一、二分别给出恢复收益出现的位置和向珍稀物种传递时的损失；问题三、四把承载条件与长期轨道联系起来；问题五再比较污染、入侵叠加后的相对状态。各模块共用时间口径、情景设定和上游输入，承担不同判断，彼此不替代。

\noindent 现有资料可以支持方向性的情景比较和模型内机制解释。若连续河段监测、标准化CPUE、珍稀物种年度调查及污染分级资料齐全，单区结果才能按河段复算，长期动力学参数也才能与实际营养级数据对应。

